\documentclass[10pt,a4paper]{amsart}
\usepackage[margin=19mm]{geometry}
\usepackage{amsmath,amssymb,mathtools}
\usepackage[scr=rsfs,cal=euler]{mathalfa}
\usepackage{xfrac}
\usepackage[unicode,hidelinks,bookmarks=false]{hyperref}
\newcommand{\ie}{i.e.\ }
\newcommand{\cf}{cf.\ }
\newcommand{\resp}{respectively\ }
\newcommand{\Subsection}{\subsection}
\providecommand{\nobreakdash}{\nobreak\hbox{-}}
\setlength{\emergencystretch}{2em}
\multlinegap0pt
\usepackage{bm}
\usepackage{bbm}
\usepackage{dsfont}
\usepackage{xspace}
\usepackage{cancel}
\usepackage{mathtools}
\usepackage{amsthm}
\makeatletter
\@ifundefined{lemma}{\newtheorem{lemma}{Lemma}}{}
\makeatother
\title[Friedman's $ \mathsf{WD} $ is not parameter-free sequential]{Friedman's $ \mathsf{WD} $ is not parameter-free sequential}
\author{Juvenal Murwanashyaka}
\date{Source: arXiv:2509.14222v1 (CC BY 4.0)}
\begin{document}
\maketitle

\noindent\emph{Source and licence.} Friedman's $ \mathsf{WD} $ is not parameter-free sequential. Version arXiv:2509.14222v1 (CC BY 4.0), distributed under the Creative Commons Attribution 4.0 International licence, as recorded in the arXiv licence metadata for this exact version and shown on the abstract page https://arxiv.org/abs/2509.14222v1. Full rights notice: licenses/arxiv-2509.14222-CC-BY-4.0.txt.\par\smallskip

\noindent\textit{Editorial scope.} Counted unit: the Lemma whose source label is \texttt{SecondLemmaGood} in the article (source0.tex lines 982--988, source label \texttt{SecondLemmaGood}), delivered together with the article's own complete original proof of it (source0.tex lines 989--1337, one \texttt{proof} environment of 349 lines). No mathematics is written, repaired or re-derived: statement and proof are the article's own text line for line. Required context retained uncounted: none. Every definition, display and label of those lines is retained unchanged. Omitted from this package and not redistributed: 1--981, 1338--1967. Nothing inside a retained excerpt is omitted or altered: every retained body is delivered exactly as the article's lines are written, so no omission receipt of any kind accompanies this package. Citations: the retained excerpts carry no citation command at all, so no bibliography entry of the article is transcribed and no reference is redistributed. Reference adaptation: none was needed and none was made; every reference inside the delivered unit resolves inside it, so no unresolved reference or citation is produced. Printed slips kept literal: no printed word, symbol, hypothesis or number of the counted statement or of its proof was altered. Layout and class: the article's own class is \texttt{memoir} with packages including inputenc, babel, amsmath, amsthm, thmtools, mathtools, amsfonts, amssymb, makeidx, graphicx, tikz-cd, float, authblk, forest; the wrapper is the lane's pinned \texttt{amsart} editorial wrapper at 10pt on A4, which restates the theorem-like environments the retained text needs and adds the article's own macro definitions that the retained text needs (none), with extra wrapper packages (bm, bbm, dsfont, xspace, cancel, mathtools). The delivered numbering is the unit's own. Assets: no retained span contains a figure environment, a graphics-inclusion command, a PDF-page inclusion, a table or a TikZ picture; no image file is copied into this package, the package declares no asset dependency and no image file is redistributed with it. Licence: version arXiv:2509.14222v1 of this article is distributed under the Creative Commons Attribution 4.0 International licence, as recorded in the arXiv licence metadata for this exact version and shown on the abstract page https://arxiv.org/abs/2509.14222v1. A full rights notice is delivered at licenses/arxiv-2509.14222-CC-BY-4.0.txt. Characters: every retained excerpt is pure ASCII, so the delivered engine, XeLaTeX with its default Latin Modern text font, reports no missing character.
 \begin{lemma}    \label{SecondLemmaGood}
 Let  $  \sigma  \in  \mathbb{P} $. 
Let  $  x_0, y_0  \in  V^{  \star }  $. 
Assume the one-element sequence $  x_0$  is  $  \mathsf{Good}_1  \left(  \sigma \right) $
and   the one-element sequence $  y_0$  is  $  \mathsf{Good}_1  \left(  \sigma \right) $. 
Then,  $  (x_0,  y_0)  $  is     $  \mathsf{Good}_2  \left(  \sigma \right) $. 
 \end{lemma}

\begin{proof}

Let  $  z  \in  \lbrace  x_0, y_0  \rbrace $. 
Since the one-element sequence $z $  is  $   \mathsf{Good}_1  \left(  \sigma  \right)  $, 
we have  $    \Delta \left(  \sigma   \,  ;   \,   z  \right)    \subseteq     \mathsf{dom} \left(  \sigma  \right)   $  
and for all $  u  \in    \Delta \left(  \sigma   \,  ;   \,   z  \right)    $  and  all finite  sets  $  A , B \subseteq   \Delta \left(  \sigma   \,  ;   \,   z  \right)   $, 
there exists  $  v  \in    \Delta \left(  \sigma   \,  ;   \,   z  \right)    $  such that  
$  \sigma (v)  =   \left(   \sigma (u)  \setminus   A   \right)  \cup   B      \,     $. 
We define the sets  $  \Delta \left(  \sigma , k  \,  ;   \,   z  \right)  $  by recursion on  $ k  \in  \omega   $: 
\begin{enumerate}
\item  $  \Delta \left(  \sigma , 0 \,  ;   \,   z  \right)   =  \lbrace   z  \rbrace $ 




\item   $  \Delta \left(  \sigma , k +1  \,  ;   \,   z  \right)   =   \Delta \left(  \sigma , k  \,  ;   \,   z  \right)   \cup  W_k  $ 
where  $  W_k $  consists of all $  v  \in  \mathsf{dom} \left(  \sigma  \right)  $  such that 
$  \sigma (v)  =    \sigma (z)  \setminus  A   $  for some  $  A  \subseteq     \Delta \left(  \sigma , k  \,  ;   \,   z \right)        \,    $.
\end{enumerate} 
We   have the following observation. 







  \begin{quote} {\bf (Claim)} \;\;\;\;\;\;    
  Let   $   z  \in    \lbrace  x_0,  y_0  \rbrace   $. 
  The following holds for all  $  k  \in   \omega   $: 
  \begin{enumerate}
  \item    $    \Delta \left(  \sigma , k   \,  ;   \,   z  \right)     \subseteq   \sigma (z)       $
  and  $  \sigma (u)  \subseteq   \sigma (z)  $  for all  $  u \in       \Delta \left(  \sigma , k   \,  ;   \,   z  \right)     \,    $. 
  
  
    
  \item   $  \sigma (s)  \setminus  \sigma (t)   \subseteq       \Delta \left(  \sigma , k   \,  ;   \,   z  \right)    $
  for all $  s, t  \in     \Delta \left(  \sigma , k   \,  ;   \,   z  \right)         \,   $. 
  
  
  
  
  
  \item  For all $  u \in   \Delta \left(  \sigma , k   \,  ;   \,   z  \right)    $  and all $  A, B  \subseteq    \Delta \left(  \sigma , k   \,  ;   \,   z  \right)   $,   there exists  $  v  \in    \Delta \left(  \sigma , k+1   \,  ;   \,   z  \right)   $
  such that    $  \sigma (v)  =   \left(   \sigma (u)  \setminus   A   \right)  \cup   B      \,     $. 

  
  

  \item    $  v   \in  \sigma (w)     $  for    all   $  w  \in    \Delta \left(  \sigma , k   \,  ;   \,   z  \right)    $  and  
all $  v  \in   \Delta \left(  \sigma , k +1  \,  ;   \,   z  \right)     \setminus    \Delta \left(  \sigma , k   \,  ;   \,   z  \right)    $.  
  \end{enumerate}
 \end{quote}  




First,  we show that (1) holds for all $ k   \in  \omega   $. 
We prove this by induction on $k$. 
We consider the base case $  k =  0 $. 
This follows from the fact that  $    \Delta \left(  \sigma , 0  \,  ;   \,   z  \right)    =  \lbrace z  \rbrace  $  
and   $  z  \in    \sigma (z)        $  by how    $  \mathbb{P} $  is defined. 
We consider the inductive case   $  k >  0 $. 
Pick  $  u  \in    \Delta \left(  \sigma , k  \,  ;   \,   z  \right)   $. 
We need to show that  $  u  \in  \sigma (z)  $  and  $  \sigma (u)   \subseteq   \sigma (z)  $. 
We have two cases: 
(i)   $  u  \in     \Delta \left(  \sigma , k-1  \,  ;   \,   z  \right)    $; 
(ii) $  u  \not\in   \Delta \left(  \sigma , k-1  \,  ;   \,   z  \right)    $. 
Case  (i) is fine since we just use the induction hypothesis. 
We consider (ii). 
Since $  u  \not\in   \Delta \left(  \sigma , k-1  \,  ;   \,   z  \right)    $,   there exists   $  A  \subseteq    \Delta \left(  \sigma , k-1  \,  ;   \,   z  \right)    $  such that  
$  \sigma (u)  =  \sigma (z)  \setminus  A  $. 
This clearly shows that  $  \sigma (u)   \subseteq    \sigma (z)   $. 
It remains to show that  $  u  \in  \sigma (z)  $. 
But, by how  $  \mathbb{P} $  is defined,  $  u  \in   \sigma (u)   \subseteq   \sigma (z)  $. 
Thus, by induction,  (1)  holds for all $ k  \in  \omega  $. 






We show that (2) holds   holds for all $ k   \in  \omega   $. 
We prove this by induction on $k$. 
The base case $k= 0 $  holds since   $   \Delta \left(  \sigma , 0  \,  ;   \,   z  \right)   $  is a singleton. 
We consider the inductive case  $  k >  0  $. 
By how  $   \Delta \left(  \sigma , k  \,  ;   \,   z  \right)   $  is defined,  there exists 
 $  v_0,  v_1,  \ldots ,  v_j  \in   \Delta \left(  \sigma , k-1  \,  ;   \,   z  \right)    $  
 and  $  A_0,  A_1,  \ldots , A_j  \subseteq    \Delta \left(  \sigma , k-1  \,  ;   \,   z  \right)  $  such that  
 \[
 \Delta \left(  \sigma , k  \,  ;   \,   z  \right)  =   \Delta \left(  \sigma , k-1  \,  ;   \,   z  \right)     \cup   
 \left\{
     \sigma^{ -1}   \left(      \sigma ( v_i )  \setminus   A_i   \right)     :    \     i  \in  \lbrace 0, 1,  \ldots  ,  j  \rbrace 
 \right\}
 \      .
 \]
Pick  $  s,  t   \in  \Delta \left(  \sigma , k  \,  ;   \,   z  \right) $.
We need to show that 
 $  \sigma (s)  \setminus  \sigma (t)   \subseteq       \Delta \left(  \sigma , k   \,  ;   \,   z  \right)    $.
 We have two cases:
  (i)  $  s  \in   \Delta \left(  \sigma , k-1  \,  ;   \,   z  \right) $; 
  (ii)  $  s  \not\in   \Delta \left(  \sigma , k-1  \,  ;   \,   z  \right) $. 
  We consider (i). 
  If  $  t  \in   \Delta \left(  \sigma , k-1  \,  ;   \,   z  \right)   $, 
  then   $  \sigma (s)  \setminus  \sigma (t)    \subseteq   \Delta \left(  \sigma , k-1  \,  ;   \,   z  \right)    \subseteq       \Delta \left(  \sigma , k   \,  ;   \,   z  \right)    $
  by the induction hypothesis. 
  Otherwise, $   \sigma (t)  =    \sigma ( v_i )  \setminus   A_i   $  for some  $  i  \in    \lbrace 0, 1,  \ldots  ,  j  \rbrace   $ 
  and hence  by the induction hypothesis 
  \[
    \sigma (s)  \setminus  \sigma (t)        \subseteq   
    \left(       \sigma (s)  \setminus \sigma ( v_i )       \right)    \cup  A_i 
    \subseteq   \Delta \left(  \sigma , k-1  \,  ;   \,   z  \right)    \subseteq       \Delta \left(  \sigma , k   \,  ;   \,   z  \right)  
    \      .
  \]
Finally,  we consider (ii). 
We have  $  \sigma (s)  =    \sigma ( v_{  \ell}   )  \setminus   A_{  \ell  }   $  for some  $  \ell  \in    \lbrace 0, 1,  \ldots  ,  j  \rbrace   $. 
If  $  t   \in    \Delta \left(  \sigma , k-1  \,  ;   \,   z  \right)   $,  then    by the induction hypothesis 
  \[
    \sigma (s)  \setminus  \sigma (t)        \subseteq        \sigma ( v_{  \ell}   )   \setminus \sigma ( t )    
    \subseteq   \Delta \left(  \sigma , k-1  \,  ;   \,   z  \right)    \subseteq       \Delta \left(  \sigma , k   \,  ;   \,   z  \right)  
    \      .
  \]
Otherwise,   $   \sigma (t)  =    \sigma ( v_i )  \setminus   A_i   $  for some  $  i  \in    \lbrace 0, 1,  \ldots  ,  j  \rbrace   $ 
  and hence  by the induction hypothesis 
  \[
    \sigma (s)  \setminus  \sigma (t)        \subseteq       
    \left(   \sigma ( v_{  \ell}   )   \setminus  \sigma ( v_i )      \right)    \cup   A_i 
    \subseteq   \Delta \left(  \sigma , k-1  \,  ;   \,   z  \right)    \subseteq       \Delta \left(  \sigma , k   \,  ;   \,   z  \right)  
    \      .
  \]  
Thus, by induction,  (2)  holds for all $ k  \in  \omega  $. 





We show that (3) holds   holds for all $ k   \in  \omega   $. 
Pick     $  u  \in  \Delta \left(  \sigma , k   \,  ;   \,   z  \right)   $  and    $  A,  B   \subseteq  \Delta \left(  \sigma , k   \,  ;   \,   z  \right)   $. 
We need to show that  there exists   $  v  \in    \Delta \left(  \sigma , k +1   \,  ;   \,   z  \right)    $  such that  
$  \sigma (v)  =      \left(   \sigma (u)  \setminus   A   \right)  \cup   B  $. 
By how  $   \Delta \left(  \sigma , k +1   \,  ;   \,   z  \right)      $  is defined,  
  we need to show that 
there exists  $  C  \subseteq    \Delta \left(  \sigma , k   \,  ;   \,   z  \right)   $  such that 
\[
 \left(   \sigma (u)  \setminus   A   \right)  \cup   B    =    \sigma (z)   \setminus  C
 \       .
\]
We have  $ \sigma  (u)  \cup    A  \cup  B   \subseteq   \sigma (z)  $  by clause (1) of the claim.
By clause (2) of the claim,    there exists    $  D  \subseteq      \Delta \left(  \sigma , k   \,  ;   \,   z  \right)   $  such that 
$  \sigma  (u)  =   \sigma (z)  \setminus  D  $. 
Hence 
\begin{align*}
 \left(   \sigma (u)  \setminus   A   \right)  \cup   B    
 &=  
 \left(   \sigma (z)  \setminus   \left(  A  \cup   D  \right)     \right)  \cup   B  
 \\
 &=  \sigma (z)   \setminus  C
&    \mbox{  where  }   C  :=    \left(  A  \cup   D  \right)     \setminus    B  
 \        .
\end{align*}
Thus,   (3)  holds  for all $  k  \in  \omega  $.  






Finally,    we   show that (4) holds  for all  $  k  \in  \omega  $.
Pick  $  w  \in    \Delta \left(  \sigma , k   \,  ;   \,   z  \right)    $   and     $  v  \in   \Delta \left(  \sigma , k +1  \,  ;   \,   z  \right)     \setminus    \Delta \left(  \sigma , k   \,  ;   \,   z  \right)    $. 
By clause (2) of the claim 
\[
v  \in   \cap   \sigma   \left[    \Delta \left(  \sigma , k   \,  ;   \,   z  \right)     \right]  
\      \      \mbox{  or   }      \         \   
v    \not\in   \cup   \sigma   \left[    \Delta \left(  \sigma , k   \,  ;   \,   z  \right)     \right]  
\      .
\]
By clause  (1) of the claim,  $  v  \in   \Delta \left(  \sigma , k +1  \,  ;   \,   z  \right)     \subseteq  \sigma (z)  $. 
Hence,  $  v  \in   \cap   \sigma   \left[    \Delta \left(  \sigma , k   \,  ;   \,   z  \right)     \right]    $
since  $  z  \in     \Delta \left(  \sigma , k   \,  ;   \,   z  \right)           \,     $. 
In particular,    $  v   \in  \sigma (w)        \,    $. 





It follows clauses (2)-(3) of the  claim   and the closure properties that define    $      \Delta \left(  \sigma   \,  ;   \,   z  \right)   $  that 
\[
  \Delta \left(  \sigma   \,  ;   \,   z  \right)    =  \bigcup_{  k  \in  \omega  }      \Delta \left(  \sigma , k  \,  ;   \,   z  \right)   
  \       . 
\]






We show that    $  (x_0,  y_0)  $  is     $  \mathsf{Good}_2  \left(  \sigma \right) $.
It suffices to construct an increasing sequence  of isomorphisms 
\[
f_k  :    \left(   \Delta \left(  \sigma , k  \,  ;   \,   x_0   \right)    \,  ,  \,   \in^{  \sigma }   \right)  
   \to  
   \left(     \Delta \left(  \sigma , k  \,  ;   \,   y_0  \right)       \,  ,  \,   \in^{  \sigma }   \right)  
\     \    \mbox{  for  }  k  \in  \omega 
\     .
\]
By setting  $  f  :   \bigcup_{  k   \in   \omega }    f_k $,  we then get an isomorphism 
\[
f  :    \left(   \Delta \left(  \sigma   \,  ;   \,   x_0   \right)    \,  ,  \,   \in^{  \sigma }   \right)  
   \to  
   \left(     \Delta \left(  \sigma   \,  ;   \,   y_0  \right)       \,  ,  \,   \in^{  \sigma }   \right)  
\]
that witnesses that    $  (x_0,  y_0)  $  is     $  \mathsf{Good}_2  \left(  \sigma \right)  $. 





We construct the sequence    $  \langle f_k  :   \   k  \in  \omega   \,  \rangle  $  by recursion. 
We have  $  f_0  \left(  x_0  \right)   = y_0 $  since 
 $    \Delta \left(  \sigma , 0  \,  ;   \,   z   \right)    =  \lbrace  z  \rbrace  $  for all $  z  \in  \lbrace x_0,  y_0  \rbrace  $. 
This map is an isomorphism since   $  x_0  \in  \sigma \left(  x_0  \right)   $  and    $  y_0  \in   \sigma \left(  y_0  \right)   $.
Assume we have defined  $  \langle f_j  :   \    j  \leq  k    \,  \rangle  $. 
We define  $  f_{ k+1}  $: 
\begin{enumerate}
\item  If   $   v  \in    \Delta \left(  \sigma , k  \,  ;   \,   x_0   \right)    $, then 
$  f_{ k+1}  \left( v  \right)    :=    f_k (v)        \,   $. 



\item   If  $  v  \not\in     \Delta \left(  \sigma , k  \,  ;   \,   x_0   \right)     $  
and   hence  $  \sigma ( v )  =   \sigma  ( x_0  )   \setminus  A_v    $
where   $   A_v     \subseteq    \Delta \left(  \sigma , k  \,  ;   \,   x_0   \right)   $, 
then 
\[
 f_{ k+1}  \left( v  \right)     
 := 
 \sigma^{  -1}  \Big(      \sigma  ( y_0 )   \setminus  f_k \left[  A_v  \right]      \Big)  
\      .
\]
\end{enumerate}
First, we show that  $ f_{ k+1}  $  is one-to-one. 
Since  $  f_k$  and $  \sigma $  are both one-to-one,  
it suffices to show that  $  f_{ k+1}  (v)  \not\in    \Delta \left(  \sigma , k  \,  ;   \,   y_0   \right)   $  
for all $  v  \in   \Delta \left(  \sigma , k+1  \,  ;   \,   x_0   \right)   \setminus   \Delta \left(  \sigma , k  \,  ;   \,   x_0   \right)   $.
Assume for the sake of  a contradiction,  there exist
$  w  \in     \Delta \left(  \sigma , k  \,  ;   \,   x_0   \right)   $  and   
$  v  \in   \Delta \left(  \sigma , k+1  \,  ;   \,   x_0   \right)   \setminus   \Delta \left(  \sigma , k  \,  ;   \,   x_0   \right)   $  
such that    
$
f_k (w)  =    f_{ k+1}  \left( v  \right)   
$.
We have 
\[
\sigma \left(   f_k \left( w  \right)     \right)    = 
\sigma \left(   f_{ k+1}  \left( v  \right)     \right)  =      \sigma  ( y_0 )   \setminus  f_k \left[  A_v  \right]   
\       .
\]
Since  $  f_k $  is an isomorphism, we must also have 
\[
\sigma (w)  =     \sigma  ( x_0 )   \setminus   A_v 
\       .
\]
But since  $  \sigma (v)  =   \sigma  ( x_0  )   \setminus  A_v     $  and  $  \sigma  $  is one-to-one, 
we get  that  $  v  =   w  $,  which contradicts the assumption that  
$  v  \not\in     \Delta \left(  \sigma , k  \,  ;   \,   x_0   \right)   $.
Thus,   $  f_{ k+1}  $  is  one-to-one. 









It remains to show that that  $f_{ k+1}  $  is an isomorphism. 
Pick    $  v, w  \in  \Delta \left(  \sigma , k+1  \,  ;   \,   x_0   \right)   $. 
We need to show that  $  v  \in   \sigma ( w)   $  if and only if  $   f_{ k+1} ( v )   \in     \sigma  \left(   f_{ k+1} ( w )    \right)     \,    $.
We have two cases: 
(I) $   w  \in  \Delta \left(  \sigma , k  \,  ;   \,   x_0   \right)    $; 
(II)   $   w    \not\in  \Delta \left(  \sigma , k  \,  ;   \,   x_0   \right)    $  and 
$  \sigma ( w )  =   \sigma  ( x_0  )   \setminus  A    $
where    $   A    \subseteq    \Delta \left(  \sigma , k  \,  ;   \,   x_0   \right)   $.
We consider case (I). 
If  $  v  \in  \Delta \left(  \sigma , k \,  ;   \,   x_0   \right)   $,  
then   $  v  \in   \sigma ( w)   $  if and only if  $   f_{ k+1} ( v )   \in     \sigma  \left(   f_{ k+1} ( w )    \right)        $
since  $ f_k $  is an isomorphism and  $ f_{ k+1} $  agrees with $  f_k  $  on   $  \Delta \left(  \sigma , k \,  ;   \,   x_0   \right)   $. 
Assume   $  v    \not\in  \Delta \left(  \sigma , k \,  ;   \,   x_0   \right)   $. 
Since  $  f_{ k+1}  $  is one-to-one,   we also have  $  f_{ k+1}  \left(  v  \right)     \not\in  \Delta \left(  \sigma , k \,  ;   \,   y_0   \right)   $. 
By clause (4) of the claim
\[
v  \in   \sigma ( w)     \     \mbox{  and  }        \   
 f_{ k+1} ( v )   \in    \sigma  \left(   f_{ k} ( w )    \right)     =  \sigma  \left(   f_{ k+1} ( w )    \right)      
 \            .
\]








We consider Case (II):   $  \sigma ( w )  =   \sigma  ( x_0  )   \setminus  A    $
and  $  \sigma   \left(   f_{ k+1}    (w )    \right)   =   \sigma  (\left(   y_0    \right)   \setminus   f_k  \left[ A  \right]      $ 
where    $   A    \subseteq    \Delta \left(  \sigma , k  \,  ;   \,   x_0   \right)   $.
We have two subcases: 
(IIa)  $  v  \in    \Delta \left(  \sigma , k \,  ;   \,   x_0   \right)    $; 
(IIb)  $  v  \not\in   \Delta \left(  \sigma , k \,  ;   \,   x_0   \right)       \,    $. 
We consider (IIa). 
Since  $f_k $  is an isomorphism and  $  f_{ k+1}  $  agrees with  $  f_k  $  on   $   \Delta \left(  \sigma , k  \,  ;   \,   x_0   \right)   $
\begin{align*}
 v  \in   \sigma ( w)       \   
 &  \Leftrightarrow     \  
 v   \in  \sigma ( x_0 )     \;   \wedge    \;  v  \not\in  A  
 \\
  &  \Leftrightarrow     \  
 f_k  (v)    \in  \sigma  \left(   y_0  \right)     \;   \wedge    \;    f_k  (v)     \not\in  f_k  \left[  A  \right] 
   \\
  &  \Leftrightarrow     \  
 f_{k+1}  (v)    \in  \sigma  \left(   y_0  \right)     \;   \wedge    \;    f_k  (v)     \not\in  f_k  \left[  A  \right] 
       \\
  &  \Leftrightarrow     \  
   f_{k+1}  (v)     \in    \sigma   \left(    f_{k+1}  (w)    \right)  
   \        .
\end{align*}




Finally,  we consider (IIb). 
Since   $  v  \not\in   \Delta \left(  \sigma , k \,  ;   \,   x_0   \right)    $,  
we also have     $  f_{ k+1}   \left( v  \right)    \not\in   \Delta \left(  \sigma , k \,  ;   \,   y_0   \right)    $. 
By clause (4)  of the claim  
\[
v  \in  \cap   \sigma  \left[       \Delta \left(  \sigma , k \,  ;   \,   x_0   \right)     \right]   
\    \     \mbox{  and   }     \      \  
f_{ k+1}  \left( v  \right)   \in  \cap   \sigma  \left[       \Delta \left(  \sigma , k \,  ;   \,   y_0   \right)     \right]   
\     .
\] 
Hence,  since   $  A  \subseteq      \Delta \left(  \sigma , k \,  ;   \,   x_0   \right)    $
\[
v  \in     \sigma  ( x_0  )   \setminus  A   =   \sigma ( w )  
\     \    \mbox{  and   }    \     \
f_{ k+1}  \left( v  \right)   \in        \sigma  (\left(   y_0      \right)   \setminus   f_k  \left[ A  \right]        =   \sigma   \left(   f_{ k+1}    (w )    \right)  
\      .
\]
This completes the proof that  $  f_{ k+1}  $  is an isomorphism.
\end{proof}

\end{document}
